

\illustration{0.3}{quadrilaterals.jpg}{
    Block puzzle with quadrilaterals.
    From the top left, in clockwise order:
    kite, trapezium, parallelogram, rhombus, rectangle, square. \\
    Modified from \href{https://www.discountschoolsupply.com/stem-curriculum/stem/mathematics/excellerations-wooden-shape-puzzles---set-of-10/p/31924}{Excellerations\textregistered{} Wooden Shape Puzzles}
}
% Font: https://www.whatfontis.com/FF_VI-Hai-Duong.font?text=parallelogram#ct

% optionally define variables/limits for this problem


For the longest time you could keep your toddlers happy by letting them play with
triangular, square, and circular wooden blocks that fit exactly through perfectly sized holes.
After letting them play a bit too long, they completely mastered this game and are
now bored, preventing you from fixing the bugs in your code.

Just now, they decided to reverse roles and started screaming planar coordinates at you,
insisting that you determine which shape each four points make:
kite, trapezium, parallelogram, rhombus, rectangle, square, or none of those.
You do not have time for this, since your bugs still need fixing.
Instead, you write a new program to answer your toddlers' questions, ideally without bugs.

The definitions for the quadrilateral shapes are as follows:
\vspace{-0.75em}
\begin{itemize}
    \setlength\itemsep{0pt}
    \item A square has four right angles and four sides with equal lengths.
    \item A rectangle has four right angles.
    \item A rhombus has four sides with equal lengths.
    \item A parallelogram has two pairs of parallel sides.
    \item A trapezium has one pair of parallel sides.
    \item A kite has reflection symmetry across a diagonal.
\end{itemize}

